\subsection{Height Formula}
\label{sec:height}

The height of an element is the sum of its content height and its vertical padding:

\begin{equation}
\label{eq:height}
h = (n \cdot 6 + 2 \cdot d \cdot w) \; U.
\end{equation}

The term $n \cdot 6\,U$ represents the intrinsic content height. Each text line occupies $6\,U$ ($= 1.5 \times \text{fontSize}$), which corresponds to the $1.5$ line-height ratio adopted to satisfy the WCAG~2.1 resilience criterion~\cite{wcag21text} and validated by readability research~\cite{rello2016}.

The term $2 \cdot d \cdot w \cdot U$ represents the total vertical padding (top and bottom combined). It grows linearly with both density $d$ and wrapping level $w$.

\subsubsection{Single-Line Controls}

For the most common case --- a single-line bounded control ($n = 1$, $w = 1$):
\begin{equation}
h = (6 + 2d) \; U.
\end{equation}

At the default density $d = 1.5$:
\begin{equation}
h = 9\,U.
\end{equation}

At $\text{fontSize} = 16\,\text{px}$ ($U = 4\,\text{px}$), this yields $h = 36\,\text{px}$.

\subsubsection{Multi-Line Blocks}

For multi-line bounded blocks ($w = 2$):
\begin{equation}
h = (6n + 6) \; U \quad \text{at } d = 1.5.
\end{equation}

This adds one full text-line height ($6\,U$) of total vertical padding, split equally above and below the content region.

\subsubsection{Height Reference}

Table~\ref{tab:height-reference} shows the resulting heights across all wrapping levels at the default density.

\begin{table}[ht]
\centering
\caption{Element height at default density $d = 1.5$, $n = 1$.}
\label{tab:height-reference}
\begin{tabular}{cccc}
\toprule
$w$ & Formula & Height ($U$) & Height (px) \\
\midrule
0 & $6\,U$ & 6 & 24 \\
1 & $9\,U$ & 9 & 36 \\
2 & $12\,U$ & 12 & 48 \\
3 & $15\,U$ & 15 & 60 \\
\bottomrule
\end{tabular}
\end{table}

